\section*{Use of LLM-Based Tools}
Tools based on large language models (LLMs) were used substantively in this work: for proof
attempts and for checks of arguments and of candidate lemmas, for exploratory programs and
searches for counterexamples, for large parts of the Lean formalization under the direction
of the author, and for drafting and revising the text. Results of computer searches were
treated as evidence, not as proofs of universal claims. The decision procedures and the
bounds on the length of witnesses do not rest on these tools: they are established by proofs
checked by the Lean kernel (\cref{sec:mechanization}). The size of the flattening, the
running time of the verification of a guessed witness, and the complexity classification are
proved on paper only. The author is responsible for all statements and proofs in the paper.
